\label{sec:resumo_atividades}

The work used the simulated events of the previous stage, whose files date from October 2025: \qty{100}{\MeV} deposits, rescaled to \qty{10}{\MeV}, in one octant of the third far detector module (FD3) of DUNE (Deep Underground Neutrino Experiment), read out by \num{6912} channels. In December 2025, a prototype graph neural network (GNN) was tested, followed, in February and March 2026, by the GATv2 and GPS encoders, with a hierarchical regression head. In June, the code was rebuilt as an installable package; the work corrected two geometry defects and three hot channels, mirrored the training events into the other octants and restricted the graphs to the lit channels, with \num{17} features per node and \num{23} per edge. In July, a simple gradient-boosting model outperformed the GNN of the time, and the XGBoost and LightGBM boosters were added as reference models.

On 4 August 2026, the boosters were tuned, and a first comparison trained \num{11} of the \num{14} encoders on a quarter of the training events. A diagnosis of the gap to the boosters led to the use of the \num{7490} training events and \num{33} event-level scalars, which summarise the event's light. In the final comparison of 8 August, the training of \num{11} matched-capacity encoders, with a maximum deviation of \qty{1.6}{\percent} from \num{1843395} parameters, totalled \qty{11.42}{\hour}.

In this report, without retraining, the boosters, tuned and trained on light with data augmentation, were re-evaluated on the same light as the encoders, and the confidence intervals and the light barycentre were computed, all on the same test events. The outputs are a paper with the results for the 2026 Scientific Initiation Meeting of the Aeronautics Institute of Technology (ENCITA), project slides and a repository of \num{155} commits and \num{172} Python files, with \num{18025} lines and \num{130} test functions.

On the \num{937} test events, the four encoders with the lowest error (two variants of GATv2 and two of GPS), indistinguishable from one another, have a mean three-dimensional error (mean distance from the predicted to the true vertex) of \num{25.95} to \qty{26.27}{\cm} and a root-mean-square error of the three axes together (combined RMSE) of \num{18.55} to \qty{18.80}{\cm}. XGBoost, on the same light, has a combined RMSE of \qty{20.95}{\cm}, but errs less on the brightest events. The network of the previous stage has \qty{20.86}{\cm}, a value computed from the published per-axis RMSE values, which come from a different partition \citep{vicentin2024study}. APEX (Aluminum Profiles with Embedded X-ARAPUCA), from a different simulation \citep{marinho2025apex}, has a combined width (root mean square of the resolution widths of the three axes, read from a plot) of \qty{38.0}{\cm} at the mean light per test event with its average light yield, \qty{31.8}{\cm} with the minimum and \qty{25.9}{\cm} at equal energy. In $z$, the long axis, against \num{20.15} to \qty{20.46}{\cm} for the four encoders, APEX has \qty{24.8}{\cm} with the average yield, above the confidence interval of each of them, and \qty{19.4}{\cm} with the minimum, within it; only at equal energy, where it detects \num{3.4} times as much light, is it better (\qty{13.7}{\cm}).
